\documentclass[12pt,a4]{report}
\usepackage{textcomp}
\usepackage{fancyhdr}
\usepackage{graphicx}
\usepackage{titletoc,tocloft}
\usepackage{enumitem}
\usepackage{amsmath,amssymb}
\evensidemargin 0.5in
\oddsidemargin 0.0in
\topmargin 0.5in
\textwidth 6.5in
\headheight 0.3in
\headsep 0in
\usepackage[]{hyperref}
\usepackage{tikz}
\usetikzlibrary{arrows,shapes,snakes,automata,backgrounds,petri}
\usepackage{tabularx}
\usepackage{pgfgantt}
\usepackage[margin=1in]{geometry}
\usepackage{lscape}
\usepackage{array}
\usepackage{longtable}
\usepackage{multirow}
\usepackage[titletoc]{appendix}
\renewcommand{\bibname}{References}
\setcounter{secnumdepth}{5}
\usepackage{float}
\usepackage{subcaption}
\usepackage{multirow}
\usepackage{longtable}
\usepackage{geometry}
\geometry{margin=1in}
\usepackage{caption}
\usepackage{tocloft}
\usepackage[ruled,vlined]{algorithm2e}
\newcommand{\cellpadding}{\renewcommand{\arraystretch}{1.2}}
\newcommand{\headerpadding}{\renewcommand{\arraystretch}{1.5}}
\newcolumntype{P}[1]{>{\raggedright\arraybackslash}p{#1}}
\textheight 8.5in
\newcommand{\pa}{\hspace{3ex}}
\newcommand{\be}{\begin{equation}}
\newcommand{\ee}{\end{equation}}
\newcommand{\fs}{\hspace{2ex}}
\newcommand{\no}{\noindent}
\newcommand{\hs}{\hspace{1ex}}
\newtheorem{myth}{Theorem}
\newcommand{\smci}{{\bf smci }}
\newcommand{\smcd}{{\bf smcd }}
\newcommand{\smc}{{\bf smc }}
\newcommand{\ab}{{\bf sc4c }}
\newcommand{\hub}{H\"ubler }
\newcommand{\ds}{D^m_f}
\newtheorem{prop}{Proposition}
\newtheorem{cor}{Corollary}
\newtheorem{defn}{Definition}
\newcommand{\rr}{{\cal R }^2 }
\newcommand{\zz}{{\cal Z }^2 }
\newcommand{\at}{{A_o}^T}
\newcommand{\ato}{{A_1}^T}
\newcommand{\ath}{{\at}_{\hat{h}}}
\newcommand{\ata}{\tilde{\at}}
\newcommand{\on}{(s_1,\hat{s_1}]}
\newcommand{\tw}{(\hat{s_1},b]}
\newcommand{\hb}{\hat{B}}
\newcommand{\hfo}{{\hat{f}}_1}
\newcommand{\xy}{(x,y)}
\newcommand{\txy}{(t_x,t_y)}
\renewcommand{\r}{{\cal R}}
\newcommand{\ct}{\hat{C}_1^T}
\newcommand{\xh}{X_{\hat{h}}}
\newcommand{\ball}{B((x_o,y_o),r)}
\newcommand{\snbd}{S((x_o,y_o),k)}
\newcommand{\xyo}{(x_o,y_o)}
\newtheorem{obs}{Observation}
\newtheorem{lemma}{Lemma}
\newcommand{\CC}{{\cal C }}
\newcommand{\DD}{{\cal D }}
\newcommand{\cset}{{\{C\}}}
\newcommand{\dset}{{\{D\}}}
\newcommand{\sdset}{{\sigma_{\dset}}}
\newcommand{\dsetm}{{m\dset}}
\newcommand{\sdsetm}{{\sigma_{\dsetm}}}
\newcommand{\od}{{O_D}}
\newcommand{\id}{{I_D}}
\newcommand{\ido}{{\id_1}}
\newcommand{\idt}{{\id_2}}
\newcommand{\odo}{{\od_1}}
\newcommand{\odt}{{\od_2}}
\newcommand{\rido}{{R(\ido)}}
\newcommand{\ridt}{{R(\idt)}}
\newcommand{\rodo}{{R(\odo)}}
\newcommand{\rodt}{{R(\odt)}}
\newcommand{\xe}{$x$-extent}
\newcommand{\ye}{$y$-extent}
\newcommand{\kxo}{k_{x1}}
\newcommand{\kxt}{k_{x2}}
\newcommand{\kxth}{k_{x3}}
\newcommand{\kyo}{k_{y1}}
\newcommand{\kyt}{k_{y2}}
\newcommand{\kyth}{k_{y3}}
\newcommand{\Ra}{{\Rightarrow}}
\newcommand{\is}{\Rightarrow \hspace{10mm}}
\newcommand{\vs}{\vspace{0.2in}}
\newcommand{\im}{{\lfloor m \rfloor}}
\newcommand{\fm}
\newcommand{\dm}{{\Delta m}}
\newcommand{\ic}{{\lfloor c \rfloor}}
\newcommand{\fc}
\newcommand{\dc}{{\Delta c}}
\newcommand{\yi}{{y_i}}
\renewcommand{\xi}{{x_i}}
\newcommand{\alpp}{{{\alpha}_{2}}}
\newcommand{\alp}{{{\alpha}_{1}}}
\newcommand{\bett}{{{\beta}_{2}}}
\newcommand{\bet}{{{\beta}_{1}}}
\newcommand{\lqr}{{\lfloor q \rfloor}}
\newcommand{\ramax}{R_{max}^A}
\newcommand{\ramin}{R_{min}^A}
\newtheorem{property}{Property}
\newcommand{\sla}{{{S_{\lambda}}}}
\newcommand{\vf}{{\vspace{0.50in}}}
\newcommand{\ms}{\;}
\renewcommand{\appendixname}{Annexure}

% Define counters for modules, submodules, and sub-submodules
\newcounter{lofmodule}
\newcounter{lofsubmodule}[lofmodule]
\newcounter{lofsubsubmodule}[lofsubmodule]

% Reset lower counters when higher-level increments
\renewcommand{\thelofmodule}{\arabic{lofmodule}}
\renewcommand{\thelofsubmodule}{\thelofmodule.\arabic{lofsubmodule}}
\renewcommand{\thelofsubsubmodule}{\thelofsubmodule.\arabic{lofsubsubmodule}}

% Command for List of Figures module (level 1)
\newcommand{\lofsection}[1]{%
  \refstepcounter{lofmodule}%
  \setcounter{lofsubmodule}{0}%
  \setcounter{lofsubsubmodule}{0}%
  \addtocontents{lof}{\protect\vspace{1em}}%
  \addtocontents{lof}{\protect\noindent \textbf{\thelofmodule.\hspace{1em} #1}\protect\par}%
  \addtocontents{lof}{\protect\vspace{0.5em}}%
}

% Command for List of Figures submodule (level 2)
\newcommand{\lofsubsection}[1]{%
  \refstepcounter{lofsubmodule}%
  \setcounter{lofsubsubmodule}{0}%
  \addtocontents{lof}{\protect\vspace{0.5em}\noindent \textbf{\thelofsubmodule.\hspace{1em} #1}\protect\par}%
}

% Command for List of Figures sub-submodule (level 3)
\newcommand{\lofsubsubsection}[1]{%
  \refstepcounter{lofsubsubmodule}%
  \addtocontents{lof}{\protect\noindent \hspace{0em}\textbf{\thelofsubsubmodule.\hspace{1em}\vspace{0.5em} #1}\protect\par}%
}
\renewcommand{\contentsname}{INDEX}

\fancypagestyle{chaptertitle}{
  \fancyhf{} % clear all
  \renewcommand{\headrulewidth}{0pt}
  \renewcommand{\footrulewidth}{0pt}
}

\makeatletter
\def\@makechapterhead#1{%
  \vspace*{\fill}%
  \begin{center}
    {\Huge\bfseries CHAPTER \thechapter\\[1.5ex]#1}
  \end{center}
  \thispagestyle{chaptertitle}% 
  \vspace*{\fill}
  \clearpage
}
\def\@makeschapterhead#1{%
  \vspace*{\fill}%
  \begin{center}
    {\Huge\bfseries #1}
  \end{center}
  \thispagestyle{chaptertitle}%
  \vspace*{\fill}
  \clearpage
}
\makeatother

\newcommand{\createGanttChart}{%
    \begin{table}[htbp]
    \centering
    \begingroup
    \renewcommand\sfdefault{phv}
    \renewcommand\mddefault{mc}
    \renewcommand\bfdefault{bc}
    
    \sffamily
    \begin{ganttchart}[
        % --- Chart Styling ---
        canvas/.append style={fill=none, draw=black!5, line width=.75pt},
        hgrid style/.style={draw=black!5, line width=.75pt},
        vgrid={*1{draw=black!5, line width=.75pt}},
        title/.style={draw=none, fill=none},
        title label font=\bfseries\footnotesize,
        title label node/.append style={below=7pt},
        include title in canvas=false,
        bar label font=\mdseries\small\color{black!70},
        bar label node/.append style={left=2cm},
        bar/.append style={draw=none, fill=black!63},
        bar incomplete/.append style={fill=barblue},
        bar progress label font=\mdseries\footnotesize\color{black!70},
        group incomplete/.append style={fill=groupblue},
        group left shift=0,
        group right shift=0,
        group height=.5,
        group peaks tip position=0,
        group label node/.append style={left=.6cm},
        group progress label font=\bfseries\small,
        x unit=0.7cm,
        title label node/.append style={xshift=0.1cm}
    ]{1}{12}
      % --- Titles ---
      \gantttitle[
        title label node/.append style={below left=7pt and -3pt}
      ]{2025}{12} \\
      \gantttitlelist[
        title label node/.append style={below left=7pt and -3pt}
      ]{{"Jan"}, {"Feb"}, {"Mar"}, {"Apr"}, {"May"}, {"Jun"}, {"Jul"}, {"Aug"}, {"Sep"}, {"Oct"}, {"Nov"}, {"Dec"}}{1} \\
      % --- Tasks ---
      \ganttgroup{Project Timeline}{1}{12} \\
      \ganttbar[name=planning]{\textbf{1.} Initial Planning and Requirement Gathering}{1}{2} \\
      \ganttbar[name=research]{\textbf{2.} Research and Technology Selection}{3}{4} \\
      \ganttbar[name=dataset]{\textbf{3.} Dataset Collection and Preprocessing}{5}{6} \\
      \ganttbar[name=prototype]{\textbf{4.} Prototype Development and PoC}{7}{8} \\
      \ganttbar[name=testing]{\textbf{5.} Testing and Validation}{9}{10} \\
      \ganttbar[name=implementation]{\textbf{6.} Final Implementation and Deployment}{11}{12}\\
      \ganttbar[name=documentation]{\textbf{7.} Documentation}{10}{12}
    \end{ganttchart}
    \endgroup
    \caption{Project Timeline}
    \label{gantt:project-timeline}
    \end{table}
}
\input{psfig.sty}

\begin{document}

\pagenumbering{gobble}
\pagestyle{empty}


\vspace*{\fill}
\begin{center}

{\large  A  PROJECT REPORT ON}\\
\vspace{0.5cm}
\vspace{0.5cm}
{\bfseries \fontsize{16}{12} \selectfont{Project Name}  \\ \vspace*{1\baselineskip}}
{\bfseries \fontsize{14}{12} \selectfont{Short Desription}  \\ \vspace*{2\baselineskip}}
{\fontsize{12}{12} \selectfont Submitted towards the
 \\Partial Fulfilment of the Requirements of} \\
\vspace{1cm}
{\bfseries \fontsize{14}{12} \selectfont BACHELOR OF TECHNOLOGY \\
\vspace*{0.5\baselineskip}
IN
\vspace*{0.5\baselineskip}
\\ARTIFICIAL INTELLIGENCE AND DATA SCIENCE\\
\vspace*{1\baselineskip}} 
\fontsize{14}{12} \selectfont \textbf{By} \\
{\large \begin{center}
\begin{tabular}{l l}
Student Name 1        & Examination Seat No:  \\
Student Name 2            & Examination Seat No:  \\
Student Name 3        & Examination Seat No:  \\
\end{tabular}
\end{center}}
{\bfseries \large Under the guidance of} \\
\vspace{0.5cm}
{\bfseries \large Prof.Guide Name} \\ \vspace{0.5cm}

\begin{figure}[!h]
\centering
  \includegraphics[width=1.5in]{KKWLOGO.v1.jpg}
\end{figure}


{\small Department of Artificial Intelligence and Data Science}\\
\fontsize{13}{13} \selectfont K. K. WAGH EDUCATION SOCIETY\textquotesingle S\\
\fontsize{13}{13} \selectfont K. K. WAGH INSTITUTE OF ENGINEERING EDUCATION AND RESEARCH\\
Hirabai Haridas Vidyanagari, Amrutdham, Panchavati, Nashik, Maharashtra 422003 \\
(An Autonomous Institute since 2022, affiliated to Savitribai Phule Pune University)\\
A. Y. 2025-2026
\vspace*{\fill}
\end{center}

\pagebreak


\begin{center}

\vspace*{3\baselineskip} 
{\Large \bf{CERTIFICATE}} \\

\end{center}

\vspace{0.4cm}
\centerline{\textit{This is to certify that the project report entitled}}
\vspace*{0.6\baselineskip} 
{\bfseries \fontsize{16}{12} \centering\selectfont{Project Name}  \\}
\vspace*{0.3\baselineskip} 
{\bfseries \fontsize{14}{12} \centering\selectfont{Short Description}  \\}
\vspace*{0.6\baselineskip} 

\centerline{\textit{being submitted by}} 
\vspace*{0.5\baselineskip} 

{\large \begin{center}
\begin{tabular}{l l}
Student Name 1        & Examination Seat No:  \\
Student Name 2            & Examination Seat No:  \\
Student Name 3        & Examination Seat No:  \\
\end{tabular}
\end{center}}

\vspace*{1.5\baselineskip} 
\noindent is a bonafide work carried out by them under the supervision of Guide Name and it is approved for the partial fulfilment of the requirement of Bachelor of Technology (Artificial Intelligence and Data Science) during academic Year 2025-2026.

\vspace*{0.8\baselineskip} 
\noindent This project report has not been earlier submitted to any other Institute or University for the award of any degree or diploma.

\vspace*{3.5\baselineskip} 

\bgroup
\def\arraystretch{0.7}
\noindent
\makebox[\textwidth][s]{%
    \begin{minipage}[t]{0.45\textwidth}
        
        \begin{flushleft}
        \textbf{Guide name} \\
        Internal Guide \\
        
        \end{flushleft}
        
    \end{minipage}%
    \hfill
    \begin{minipage}[t]{0.45\textwidth}
        
        \begin{flushright}
        \textbf{Prof. Dr. S. M. Kamalapur} \\
        Head of Department, \\
        Department of AI\&DS \\
        
        \end{flushright}
        
    \end{minipage}
}
\egroup


\vspace*{2.5\baselineskip} % Reduced from 3

\begin{flushleft}

\textbf{External Examiner}\\
\vspace*{2.5\baselineskip}
\textbf{Date: }DD-MM-YYYY\\
\textbf{Place: }Nashik
\end{flushleft}

\vspace*{1.5\baselineskip}


\pagebreak

\pagenumbering{Roman}
\setcounter{page}{1}

% Define a page style for roman numbered pages with centered footer
\fancypagestyle{romanpages}{
  \fancyhf{} % clear all
  \fancyfoot[C]{\thepage} % centered page number in footer
  \renewcommand{\headrulewidth}{0pt}
  \renewcommand{\footrulewidth}{0pt}
}

% Page I: Abstract
\pagestyle{romanpages}
\begin{center}
\bfseries \fontsize{24}{12} \selectfont Abstract
\end{center}
\vspace*{1\baselineskip}
Abstract content
Keywords:
\pagebreak

% Page II: Acknowledgement
\begin{center}
\bfseries \fontsize{24}{12} \selectfont Acknowledgement
\end{center}
\vspace*{1\baselineskip}
\hspace*{0.5in}We would like to thank our project guide, Guide Name, for her guidance and support, and we will forever remain grateful for her constant encouragement extended in making this project report. Through our many discussions, she helped us form and solidify ideas, and the invaluable inputs from her and our team colleagues, along with the penetrating questions and motivation, have all contributed to the successful development of the concepts presented herein. We also extend our sincere gratitude to our Departmental Co-ordinator of Artificial Intelligence and Data Science, Dr. D. V. Medhane, and our esteemed Director, Dr. K. N. Nandurkar, for their continued support and leadership. Furthermore, we thank our \textit{sponsorship agency(include only if sponsored project)} for their generous assistance, which greatly facilitated the completion of this work.\\

\hspace*{0.3in}We would also like to thank our wonderful colleagues, parents and friends for listening our ideas, asking questions  and providing feedback and suggestions for improving our ideas.\\

\vspace*{3\baselineskip} 
\begin{flushright}
Student 1 \\
Student 2 \\
Student 3 \\
Student 4 \\
\end{flushright}

\bibliographystyle{plain}
\newpage

% Page III: INDEX (Table of Contents) - no header/footer on any page
\pagestyle{romanpages}
\tableofcontents
\addtocontents{toc}{\protect\thispagestyle{romanpages}}

\let\origaddvspace\addvspace
\renewcommand{\addvspace}[1]{}

% Page IV: List of Tables - no header/footer
\clearpage
\pagestyle{romanpages}
	\setcounter{lotdepth}{2} 
\listoftables
\addtocontents{lot}{\protect\thispagestyle{romanpages}}

% Page V: List of Figures
\clearpage
\pagestyle{romanpages}
	% \addcontentsline{toc}{section}{\bf LIST OF FIGURES}
	\setcounter{lofdepth}{2} 
\listoffigures
\addtocontents{lof}{\protect\thispagestyle{romanpages}}


\clearpage
\pagenumbering{arabic}
\setcounter{page}{1}

\fancyhf{}
\renewcommand{\footrulewidth}{0.4pt}
\lfoot{KKWIEER, Department of Artificial Intelligence and Data Science, 2025}
\rfoot{Page \thepage}

\fancypagestyle{plain}{ 
    \fancyhf{} 
    \lfoot{KKWIEER, Department of Artificial Intelligence and Data Science}
    \rfoot{Page \thepage}
    \renewcommand{\headrulewidth}{0.4pt} 
    \renewcommand{\footrulewidth}{0.4pt} 
}
\pagestyle{fancy}

\setlength{\parindent}{11mm}




\chapter{INTRODUCTION}

\section{Project Idea}

\section{Motivation of the Project}


\chapter{LITERATURE SURVEY}


\chapter{PROBLEM DEFINITION \vspace*{0.5\baselineskip} AND SCOPE}
\section{Problem Statement}


\section{Methodology}

\section{Outcome}



\chapter{PROJECT PLAN}

\section{Project Timeline}
    \createGanttChart
 
\section{Team Organization}

 
\chapter{SOFTWARE REQUIREMENT \vspace*{0.5\baselineskip} SPECIFICATION}
\section{Functional Requirements}

\section{Non-Functional Requirements}

\section{Constraints}

\section{Hardware Requirements}

\section{Software Requirements}

\section{Interfaces}

\chapter{DETAILED DESIGN}

\section{Architectural Design}

\vspace{0.3cm} 

% \lofsection{System Architecure}
% \begin{figure}[!htbp]
% \centering
% % Rotate 90 degrees to make it portrait
% \rotatebox{0}{%
%     \includegraphics[width=\linewidth, height=23cm, keepaspectratio]{./Architecture_Portrait.png}%
% }
% \caption{PolicyVault Nexus - 3-Tier Model}
% \label{architecture}
% \end{figure}


\section{Overview of Project Modules}

\section{User Interface}

\section{Data design}
\subsection{ Data structure}

\subsection{Database Description}

\section{UML Diagrams}
\vspace{0.1cm}
% \lofsection{UML Diagrams}
% \begin{itemize}
%     \item \textbf{Use Case Diagram}
%     \begin{figure}[H]
%     \centering
%     % Keep rotation 0 for portrait alignment
%     \rotatebox{0}{%
%         \includegraphics[width=18cm, height=15cm, keepaspectratio]{./usecase_diagram.png}%
%     }
%     \caption{Use Case Diagram}
%     \label{fig:usecase}
% \end{figure}
% \vspace{0.5cm}
% \end{itemize}

\section{Tools and Technology Used}


\section{Algorithm}

 \chapter{RESULTS AND DISCUSSION}
 \section{Experimental Setup}
 \subsection{Dataset}

\subsection{Performance Parameters}

\subsection{Efficiency Issues}

 \section{Software Testing}
  \subsection{Test Cases and Test Results}
 
 \section{Summary}
 \subsection{Result Analysis and Discussion}
 

 
 \chapter{CONCLUSION AND \vspace*{0.5\baselineskip} FUTURE WORK}

 \section{Conclusions}
 
 \section{Future Work}

 \section{Application}
 
\clearpage
\pagenumbering{gobble}
\pagestyle{empty}

\makeatletter
\renewcommand{\@chapapp}{Annexure}
\renewcommand{\thechapter}{\Alph{chapter}}
\makeatother

% Custom annexure chapter style
\makeatletter
\renewcommand{\@makechapterhead}[1]{%
  \vspace*{\fill}%
  \begin{center}
    {\Huge\bfseries Annexure \thechapter\\[1.5ex]#1}
  \end{center}
  \thispagestyle{chaptertitle}% disable headers/footers
  \vspace*{\fill}
  \clearpage
}
\makeatother

\begin{appendices}
\chapter{References}
% Remove any automatic bibliography title and formatting
\makeatletter
\renewenvironment{thebibliography}[1]
     {\list{\@biblabel{\@arabic\c@enumiv}}%
           {\settowidth\labelwidth{\@biblabel{#1}}%
            \leftmargin\labelwidth
            \advance\leftmargin\labelsep
            \@openbib@code
            \usecounter{enumiv}%
            \let\p@enumiv\@empty
            \renewcommand\theenumiv{\@arabic\c@enumiv}}%
      \sloppy
      \clubpenalty4000
      \@clubpenalty \clubpenalty
      \widowpenalty4000%
      \sfcode`\.\@m}
     {\def\@noitemerr
       {\@latex@warning{Empty `thebibliography' environment}}%
      \endlist}
\makeatother
\begin{thebibliography}{99}
% Chicago or Hardward style references
\bibitem{bufalino2025}
Bufalino, Jacopo, Jose Luis Martin-Navarro, Mario Di Francesco, and Tuomas Aura. "Inside Job: Defending Kubernetes Clusters Against Network Misconfigurations." \textit{Proceedings of the ACM on Networking} 3, no. CoNEXT3 (2025): 1–25.


\end{thebibliography}

\chapter{Plagiarism Report }
\chapter{Paper Published (if any)}
\chapter{Sponsorship detail (if any)}
\end{appendices}
\end{document}